% =====================================================================
%  math.tex —— 全局数学宏
%  规则：正文中的物理量一律走本文件的宏，禁止裸写下标；
%        符号表（frontmatter/notation.tex）与本文件必须一一对应。
% =====================================================================

% ---- 正体算子 --------------------------------------------------------
\DeclareMathOperator{\diag}{diag}
\DeclareMathOperator{\tridiag}{tridiag}
\DeclareMathOperator{\rank}{rank}
\DeclareMathOperator{\spann}{span}
\DeclareMathOperator*{\argmin}{arg\,min}
\DeclareMathOperator*{\argmax}{arg\,max}

% ---- 微分 ------------------------------------------------------------
\newcommand{\dd}{\mathrm{d}}                       % 直立的 d
\newcommand{\der}[2]{\frac{\dd #1}{\dd #2}}        % 全导数
\newcommand{\pder}[2]{\frac{\partial #1}{\partial #2}}

% ---- 常数与数 --------------------------------------------------------
\newcommand{\ii}{\mathrm{i}}                       % 虚数单位
\newcommand{\ee}{\mathrm{e}}                       % 自然对数底
\newcommand{\kb}{k_{\mathrm{B}}}                   % 玻尔兹曼常数
\newcommand{\kT}{k_{\mathrm{B}}T}                  % 热能
\newcommand{\epsO}{\varepsilon_{0}}                % 真空介电常数
\newcommand{\epsr}{\varepsilon_{\mathrm{r}}}       % 相对介电常数
%  面的等效相对介电常数。单独定义一个宏，因为它带上标：写成
%  \epsr^{\mathrm{face}}_i 会因 \epsr 自带下标而触发 "Double subscript"。
\newcommand{\epsface}{\varepsilon^{\mathrm{face}}}

%  ⚠ 网格节点索引必须走下面这三个宏（逗号下标形式）。
%  \Ec 与 \epsr 本身已经带下标，再写 \Ec_i 就是"同一个基上的两个下标"，
%  TeX 会报 Double subscript。这是本文件里最容易重犯的一类错误，
%  新增带角标的物理量宏时请一并在此登记。
\newcommand{\Eci}[1]{E_{\mathrm{c},#1}}                    % E_{c,i}
\newcommand{\Ecifl}[1]{E_{\mathrm{c},#1}^{\mathrm{full}}}  % E_{c,i}^{full}
\newcommand{\epsri}[1]{\varepsilon_{\mathrm{r},#1}}        % eps_{r,i}
\newcommand{\hb}{\hbar}                            % 约化普朗克常数

% ---- 物理量 ----------------------------------------------------------
\newcommand{\Ec}{E_{\mathrm{c}}}                   % 导带底
\newcommand{\Ev}{E_{\mathrm{v}}}                   % 价带顶
\newcommand{\Ef}{E_{\mathrm{F}}}                   % 费米能级
\newcommand{\Ei}{E_{i}}                            % 第 i 个亚带能级
\newcommand{\Ej}{E_{j}}                            % 第 j 个亚带能级（微扰求和用）
\newcommand{\ns}{n_{\mathrm{s}}}                   % 二维面密度
\newcommand{\nsu}{n_{\mathrm{s}}^{\max}}           % 面密度上界 sigma_eff/e
\newcommand{\meff}{m^{*}}                          % 有效质量
\newcommand{\mO}{m_{0}}                            % 自由电子质量
\newcommand{\phib}{\phi_{\mathrm{b}}}              % 肖特基势垒高度
\newcommand{\dEc}{\Delta E_{\mathrm{c}}}           % 导带带阶
\newcommand{\sigpol}{\sigma_{\mathrm{pol}}}        % 极化面电荷
\newcommand{\sigeff}{\sigma_{\mathrm{eff}}}        % 屏蔽后有效面电荷
\newcommand{\Jnn}{J_{n}}                           % 电子电流密度
\newcommand{\mumu}{\mu}                            % 迁移率
\newcommand{\NC}{N_{\mathrm{C}}}                   % 导带有效态密度
%  带有效质量修正与节点索引的导带态密度。同样必须整体定义：
%  写成 \NC\,_{\text{eff},i} 会因 \NC 自带下标而触发 Double subscript。
\newcommand{\NCeff}[1]{N_{\mathrm{C,eff},#1}}      % N_{C,eff,i}

% ---- 记号算子 --------------------------------------------------------
\newcommand{\avg}[1]{\left\langle #1 \right\rangle}   % 期望值
\newcommand{\abs}[1]{\left\lvert #1 \right\rvert}     % 绝对值
\newcommand{\norm}[1]{\left\lVert #1 \right\rVert}    % 范数
\newcommand{\T}{^{\mathsf{T}}}                        % 转置
\newcommand{\dg}{^{\dagger}}                          % Moore-Penrose 伪逆
\newcommand{\mat}[1]{\bm{#1}}                         % 矩阵

% ---- 单位 ------------------------------------------------------------
%  用法：\U{cm^{-2}}  ->  直立罗马体的单位
%  用 \ensuremath 包裹，使其在文本模式与数学模式中都能用
%  （表格与符号表里大量按文本模式书写，必须两种模式通吃）。
\newcommand{\U}[1]{\ensuremath{\,\mathrm{#1}}}
